% Infeasible Intent Handling and Fallback Policies
% For inclusion in paper Section III-D or IV (Discussion)

\subsection{Handling Infeasible, Ambiguous, and Edge-Case Intents}
\label{sec:infeasible_handling}

Real-world operator intents are not always well-formed or physically
realizable. Our system addresses three categories of problematic intents
through complementary mechanisms.

\subsubsection{Infeasible Intent Detection}

An intent is \textit{infeasible} when its constraint program is
syntactically valid but physically unrealizable---e.g., demanding
sub-millisecond latency across intercontinental paths, or routing through
a region after disabling all nodes in that region. Our 8-pass validator
detects three classes of infeasibility:

\begin{itemize}
\item \textbf{Structural infeasibility} (100\% detection): missing fields,
out-of-range entity IDs, type mismatches, and values outside physical
bounds (e.g., latency $< 2.0$\,ms). These are caught by passes 1--4
(schema, entity grounding, type safety, value ranges).

\item \textbf{Topological infeasibility} (100\% detection): constraints
that partition the network, eliminate viable paths, or cause severe
capacity loss. Pass~5 rejects contradictory constraints (e.g., disabling
a node while routing through it). Pass~7 warns when $\geq 50\%$ of nodes are disabled (severe capacity
loss) and escalates to a stronger warning at $\geq 75\%$. These
capacity thresholds are heuristic warnings outside the soundness
path---they do not block acceptance.

\item \textbf{Routing infeasibility} (100\% detection within certified
fragments): constraint combinations that are individually valid but
jointly unsatisfiable on the physical topology. Pass~8 (feasibility
certification) constructs routing witnesses using fragment-specific
algorithms (BFS, Dijkstra, hop-layered Dijkstra, Edmonds-Karp max-flow)
and rejects programs where no witness exists.
\end{itemize}

\noindent Our confusion matrix (Table~\ref{tab:confusion_matrix})
confirms the effectiveness of the 8-pass pipeline: 0\% unsafe acceptance
across all categories. Of the 30 benchmark-infeasible intents, 22 receive
\decision{Reject} and 8 receive \decision{Abstain}---none are accepted.
Pass~8 additionally identifies 32 programs among the structurally-valid
set whose latency or hop constraints cannot be satisfied on the physical
constellation topology (e.g., 30\,ms S\~{a}o Paulo--New York when the
minimum-hop path requires $\sim$63\,ms); 29/32 are independently
confirmed via a separate Dijkstra oracle.

Our adversarial safety evaluation (15 tests across resource exhaustion,
semantic conflicts, and boundary exploitation) achieves 100\% detection
(15/15), including near-total capacity removal (disabling 19/20 planes),
cross-constraint contradictions, and boundary value exploitation.

\subsubsection{Fallback Policies}

The ConstraintProgram IR includes a \texttt{fallback\_policy} field that
governs system behavior when hard constraints cannot be satisfied at
routing time:

\begin{enumerate}
\item \texttt{reject\_if\_hard\_infeasible} (default): the routing layer
refuses to compute paths and returns an explicit failure to the operator.
This is the safest option for critical intents where partial compliance
is unacceptable.

\item \texttt{relax\_soft\_first}: soft constraints are progressively
relaxed (in order of increasing penalty weight) until a feasible routing
exists. Hard constraints are never relaxed. This enables graceful
degradation for intents where approximate compliance is preferable to
total failure.

\item \texttt{report\_unsat\_core}: the system identifies the minimal
subset of constraints that cause infeasibility and reports them to the
operator, enabling informed manual intervention. This supports
diagnostic workflows where understanding \textit{why} an intent fails
is as important as resolving it.
\end{enumerate}

\noindent In our 240-intent benchmark, all programs use the default
\texttt{reject\_if\_hard\_infeasible} policy. The fallback mechanism
is designed for operational deployment where operator interaction is
available; evaluating its effectiveness under dynamic network conditions
is left to future work.

\subsubsection{Ambiguous Intent Resolution}

Ambiguous intents admit multiple valid interpretations. Our OOD
evaluation includes 5 deliberately ambiguous intents (e.g., ``optimize
the network for best performance'') to assess compiler behavior.
Key observations:

\begin{itemize}
\item The LLM compiler produces \textit{reasonable} constraint programs
for all 5 ambiguous intents (qualitative assessment), typically selecting
conservative interpretations that map to load-balancing or latency
optimization.

\item Ambiguous intents are excluded from quantitative scoring because
no unique ground truth exists. We report them separately as qualitative
evidence of graceful degradation.

\item The compiler's 6-shot prompt includes examples that implicitly
demonstrate disambiguation strategies (e.g., mapping vague performance
requests to specific constraint types), providing soft guidance without
explicit disambiguation rules.
\end{itemize}

\subsubsection{Limitations and Future Directions}

The feasibility certifier covers five constraint fragments (F1--F5)
that span the most common constraint combinations in our benchmark.
Three directions could extend coverage further:

\begin{enumerate}
\item \textbf{Extended fragment coverage}: constraints involving
\texttt{min\_cap\_reserve} or combinations of $k$-disjoint paths
with latency/hop bounds currently produce \decision{Abstain}. Adding
fragments for these combinations (e.g., via constrained max-flow)
would reduce the abstain rate.

\item \textbf{Multi-constellation generalization}: cross-constellation
evaluation (Table~\ref{tab:cross_constellation}) shows the GNN
generalizes to altitude changes but not inclination changes.
The compile--verify--ground pipeline is constellation-agnostic,
but the GNN requires retraining per orbital geometry. Extending
to heterogeneous multi-shell constellations (e.g., Starlink Gen2)
remains future work.

\item \textbf{Intent confirmation loop}: presenting the grounded
constraint program back to the operator in natural language for
confirmation before execution, closing the semantic loop between
intent and realization.
\end{enumerate}
